% ============================================================
% ch13.tex —— 第 13 章 无穷风光：0.999… 到底等不等于 1
% ============================================================
\chapter{无穷风光：\texorpdfstring{$0.999\cdots$}{0.999...} 到底等不等于 \texorpdfstring{$1$}{1}}

篇三开篇，先解决一桩你们班级群里肯定吵过的公案，然后用它撬开一扇大门。门后是本书剩余部分的主发动机：\textbf{无穷}。

\begin{kaiwei}
$0.999999\cdots$（小数点后无穷个 $9$）这个数，和 $1$ 比谁大？
\end{kaiwei}

先押个注再往下读。经验统计：大约一半人凭直觉说"$0.999\cdots$ 小一点——它永远是差一点的意思"，另一半人说"相等，但说不清为什么"。两边都觉得对方顽固。这道题的奇妙之处在于：\textbf{吵不出结果不是因为题难，而是因为双方对"$0.999\cdots$"是什么还没达成共识}。"无穷个 9 的小数"究竟是一个数、一个过程、还是一种修辞？把它说清楚，架就吵完了——而"把它说清楚"这件事，恰好就是微积分的起点。

\section{三种论证}

\paragraph{论证一（代数法，最快）。}设 $x = 0.999\cdots$。两边同乘 $10$：小数点右移一位，$10x = 9.999\cdots$。两式相减，小数部分\textbf{精确抵消}：
\[
10x - x = 9.999\cdots - 0.999\cdots = 9 \quad\Longrightarrow\quad 9x = 9 \quad\Longrightarrow\quad x = 1 .
\]
说服力强，但严格的读者会嘀咕两件事：无穷小数能"整体乘 $10$"吗？小数部分真的能"对齐相减"吗？（本章烧脑时刻专门审这两条。）先继续。

\paragraph{论证二（分数法，最朴素）。}人人都认 $\tfrac13 = 0.333\cdots$（除法竖式：$3\mid 1.000$ 商 $0.333\cdots$，永远余 $1$）。两边乘 $3$：
\[
1 = 3\times\tfrac13 = 3\times 0.333\cdots = 0.999\cdots .
\]
如果你信 $\tfrac13 = 0.333\cdots$（信，因为这是除法算出来的），那 $0.999\cdots = 1$ 就是同一句话的另一种写法。

\paragraph{论证三（无敌法，最深刻）。}回想实数的一个基本性质（数轴的"稠密"）：\textbf{两个不同的数之间，必夹着第三个数}。$0.5$ 和 $0.6$ 之间有 $0.55$；$0.999$ 和 $1$ 之间有 $0.9995$。现在假设 $0.999\cdots \ne 1$，问：哪个数比 $0.999\cdots$ 大、又比 $1$ 小？
\begin{itemize}
\item 候选 $0.9999$（四个 $9$）？它比 $0.999\cdots$ \textbf{小}（小数点后第五位它是 $0$，人家是 $9$）——不够格。
\item 候选 $0.99999$（五个 $9$）？同理不够格。
\item 任何有限个 $9$ 的候选，都被 $0.999\cdots$ 压过；任何"另一个无穷小数"候选，要么等于 $0.999\cdots$ 本身，要么首位不同就冲出了 $(0.999\cdots, 1)$ 这条缝。
\end{itemize}
\textbf{缝里塞不进任何数}。而两个不同的实数之间必有缝可塞——矛盾。所以 $0.999\cdots = 1$。

三票全过。现在做这道题真正有价值的部分——\textbf{理解它为什么违反直觉}。直觉说"$0.999\cdots$ 差一点点"，错在哪？错在把两个东西混成了一个：
\begin{itemize}
\item \textbf{一串路标}：$0.9,\ 0.99,\ 0.999,\ 0.9999, \dots$——这串数每个都\textbf{小于} $1$，永远在逼近，永远不抵达；
\item \textbf{一个目的地}："$0.999\cdots$"这个记号所指的\textbf{数}——它不是路标串里的任何一个，它是这串路标的\textbf{终点}，而终点恰好是 $1$。
\end{itemize}
直觉的"差一点点"描述的是路标（每个路标确实差），论证三说的是终点（终点不差）。\textbf{路标不是目的地}——这句话是整个微积分的第一课，请把它刻在墙上。至于"终点凭什么存在、凭什么恰好是 $1$"，需要发明一台精密的概念机器，第 14 章专造（名字叫"极限"）。

顺带拆掉一个流传很广的伪证："$0.999\cdots$ 与 $1$ 的差是 $0.000\cdots1$，一个无穷小的正数，所以差一点点。"伪在最后半句："$0.000\cdots1$"这个写法\textbf{不合法}——小数点后无穷个 $0$ 之后再写 $1$，意味着"无穷之后再有一个位置"，而无穷的意思恰恰是\textbf{没有之后}。（数学里确实有"无穷小"的严肃理论——超实数——但那里的无穷小不是这个写法，结论也不改变 $0.999\cdots = 1$。）

\section{循环小数的身份证}

\begin{kaiwei}
$0.727272\cdots$（$72$ 循环）是哪个分数？
\end{kaiwei}

老三样：设 $x = 0.7272\cdots$；两边乘 $100$（左移两位，恰好移过一个循环节）：$100x = 72.7272\cdots$；相减（循环部分精确抵消）：$99x = 72$，解得
\[
x = \frac{72}{99} = \frac{8}{11} .
\]
（顺带一个彩蛋：所有分母为 $99$ 的分数自动是两位循环——$\tfrac{72}{99}$ 的循环节正是分子补零成两位"$72$"。同理 $\tfrac1{7} = 0.\overline{142857}$ 是六位循环，因为 $7 \mid 999999$。）

现在反过来问一个更深的问题：\textbf{为什么每个循环小数都是分数？为什么每个分数都循环？}第 1 章的余数早有伏笔，把它说完。

\textbf{分数 $\Rightarrow$ 循环}。做除法 $a\div b$ 的竖式时，每一步的关键状态只有一个：\textbf{当前的余数}（$0$ 到 $b-1$ 之间的一个整数）。余数只有 $b$ 种身份（$b$ 个房间，第 1 章的分房术）。竖式每商一位，余数换一次房——房间有限，\textbf{最多 $b$ 步之内必然撞房}（抽屉原理）。撞房那一刻，后面的商开始逐字重复：循环。若余数撞上 $0$，除尽了，是有限小数（可看成"循环 $0$"）。结论：\textbf{分数的十进制展开要么有限、要么循环，且循环节长度小于分母}。

\textbf{循环 $\Rightarrow$ 分数}。刚才的"乘 $10^k$ 相减"操作就是证明：任何循环小数都能这样解出分数。

\begin{dingli}
分数 $\iff$ 有限或循环小数。两个世界一一对应（把有限小数看成"$0$ 循环"，并注意 $0.999\cdots$ 式的重复表示：$0.4999\cdots = 0.5$）。
\end{dingli}

推论震撼：$\pi = 3.14159265\cdots$ 的展开\textbf{不循环}（这是定理，不是猜想——兰伯特 1761 年证明 $\pi$ 是无理数）。由本节定理的逆否命题，\textbf{$\pi$ 不是分数}。人类花到 1761 年才严格知道：连"圆周长比直径"这么基本的量，都挤不进分数的世界。数的版图从这一刻开始扩张——分数之外，另有江山。

\section{芝诺的乌龟：无穷多个数加起来}

无穷的概念一旦立起来，立刻能干一件古人认为自相矛盾的事：\textbf{把无穷多个数加成一个确定的数}。

古希腊人芝诺（约公元前 450 年）出了一道两千五百年的难题——\textbf{阿喀琉斯追乌龟}。阿喀琉斯（希腊第一飞毛腿）速度是乌龟的十倍，让乌龟先爬 $100$ 米。裁判发令：
\begin{itemize}
\item 等阿喀琉斯跑完 $100$ 米（到达乌龟出发点），乌龟又爬了 $10$ 米；
\item 等他跑完这 $10$ 米，乌龟又爬了 $1$ 米；
\item 跑完 $1$ 米，乌龟又爬 $0.1$ 米……
\end{itemize}
芝诺说：如此追法，\textbf{永远有下一段}，所以阿喀琉斯永远追不上乌龟。这个结论显然错（现实里随便追），但两千年来"错在哪"说不清——因为我们和芝诺一样，不知道怎么处理"无穷多个数相加"。

把要跑的距离老老实实加起来：
\[
100 + 10 + 1 + 0.1 + 0.01 + \cdots \;=\; ?
\]
无穷多个正数相加——总和会不会也是无穷？硬着头皮按规则算。设总和为 $S$。观察结构：\textbf{从第二项起的部分，恰好是整体 $S$ 的十分之一}（每项都是前项的 $\tfrac1{10}$，整串缩小十倍）：
\[
S = 100 + \frac{S}{10} \quad\Longrightarrow\quad \frac{9S}{10} = 100 \quad\Longrightarrow\quad S = \frac{1000}{9} = 111.111\cdots .
\]
\textbf{无穷多个数，加出了一个有限的、确切的答案：$111\tfrac19$ 米。}阿喀琉斯跑完 $111\tfrac19$ 米就追上乌龟（之后一路领先）。

芝诺的"悖论"就此解除。他的诡计是\textbf{把有限的路程切成无穷多段，再宣称"无穷多段 $=$ 走不完"}。而计算表明：无穷多段的路程，总长可以是有限的——\textbf{无穷的过程，有限的结果}，一点矛盾都没有。芝诺缺的不是逻辑，是工具；两千五百年后我们补上工具，他的难题就变成了习题。

\begin{faming}
\textbf{土名}：无穷多个数连加 $\to$ \textbf{无穷级数}（infinite series）；连加出来的结果叫级数的\textbf{和}；加得出确定答案的叫\textbf{收敛}（converge），加不出（部分和无限膨胀或摇摆不停）的叫\textbf{发散}（diverge）。判断收敛还是发散，是接下来几章的基本功——"面对无穷连加，第一件事永远是问：它到底收敛吗？"
\end{faming}

\section{收敛不是理所当然}

上面两个级数都收敛（$0.9 + 0.09 + \cdots = 1$；$100+10+\cdots = \tfrac{1000}9$）。是不是无穷连加\textbf{总能}加拢？三个反例站出来:

\textbf{反例一}：$1 + 1 + 1 + \cdots$——部分和 $1, 2, 3, \dots$ 一路膨胀，发散。没悬念。

\textbf{反例二（最阴险）}：
\[
1 + \frac12 + \frac13 + \frac14 + \frac15 + \cdots \;=\; ?
\]
每一项都越来越小、趋近 $0$——直觉高呼"肯定收敛"。算部分和：前 $10$ 项 $\approx 2.93$，前 $100$ 项 $\approx 5.19$，前 $1000$ 项 $\approx 7.49$，前 $10^6$ 项 $\approx 14.39$……爬得极慢，但\textbf{从不减速到停止}。它究竟是收敛还是发散？答案是\textbf{发散}——而且发散的证明只有五行，与素数分布（第 22 章）直接挂钩，我们郑重地把它留成第 21 章的主菜。\textbf{先把直觉的教训记住：项趋近零，不足以保证收敛}；真正起作用的是项趋近零的\textbf{速度}。

\textbf{反例三（摇摆型）}：$1 - 1 + 1 - 1 + \cdots$——部分和 $1, 0, 1, 0, \dots$ 永远摆动，不收敛（它的"和"历史上吵了几百年，最后数学家裁决：这个级数\textbf{没有和}，问题终结）。

\begin{dongshou}
\begin{itemize}
  \yisi 把 $0.454545\cdots$ 和 $0.138888\cdots$（$8$ 循环）写成分数。（第一个：$x = 0.\overline{45}$，$100x - x = 45$，$x = \tfrac{45}{99} = \tfrac{5}{11}$。第二个：设 $x = 0.13888\cdots$，$1000x = 138.888\cdots$，$100x = 13.888\cdots$，相减 $900x = 125$，$x = \tfrac{5}{36}$。注意"错位相减"要移到循环节对齐的位置。）
  \yier 用"设 $S$ 找规律解方程"求 $1 + \tfrac13 + \tfrac19 + \tfrac1{27} + \cdots$ 的和。（$S = 1 + \tfrac S3$，$S = \tfrac32$。）
  \yier 证明芝诺变体：速度 $n$ 倍、让先 $d$ 米，则无穷段距离 $d + \tfrac dn + \tfrac d{n^2} + \cdots = \dfrac{nd}{n-1}$，恰为真实追上点。（解 $S = d + \tfrac Sn$。再用"快者速度 $v$、慢者 $\tfrac vn$"直接列方程 $vt = d + \tfrac vn t$ 解出 $t$，换算距离，两种方法对账。）
  \yisi 不计算，判断级数 $\tfrac12 + \tfrac14 + \tfrac18 + \cdots$（几何级数，公比 $\tfrac12$）与 $1 + 2 + 4 + 8 + \cdots$（公比 $2$）哪个收敛，并说明理由（部分和的趋势）。
\end{itemize}
\end{dongshou}

\begin{shaonao}
论证一（代数法）偷偷用了两条深水区规则："\textbf{无穷小数可以整体乘 $10$（移位）}"和"\textbf{无穷小数可以逐位相减}"。请审这两条：

(1) 移位合法性：$x_n = 0.\underbrace{99\cdots9}_{n}$（$n$ 个 $9$），则 $10 x_n - x_n = 9\times 0.\underbrace{0\cdots01}_{}$？不对——请精确算出 $10x_n - x_n = 9x_n$ 的值（$9 \times 0.99\cdots9$（$n$ 位）$= 8.99\cdots91$？动笔！）并观察当 $n$ 增大时它逼近谁。

(2) 相减合法性：核心是"两个\textbf{收敛}的无穷量之差 $=$ 差的收敛值"——即"极限运算与减法可交换次序"。这条规则的正式版第 14 章颁发（极限的四则运算法则）。

做完这两问你会发现：论证一不是错的，只是\textbf{预支了第 14 章的信用}。而论证三不预支任何信用——这就是它"最深刻"的含义：三种论证不是三份重复的证明，是三种深度的证明。
\end{shaonao}

\begin{chengshi}
"循环小数 $=$ 分数"我们双向说清了（余数抽屉 + 错位相减），但对"错位相减对无穷位合法"的严格背书（收敛级数的线性性）放在第 14 章。"$\pi$ 无理"的证明远超本书工具（兰伯特的连分数证明是大学内容）。芝诺悖论其实有四条（二分法、阿喀琉斯、飞矢、运动场），本章只审了最出名的一条；"飞矢不动"将在第 15 章的烧脑时刻被导数当场击落——芝诺会在这本书里被连续传唤两次。
\end{chengshi}

$0.999\cdots = 1$ 的三种论证里，最深刻的是论证三——"找不到夹缝里的数"。下一章，刘徽用同一把刀，把这句话锻造成整个微积分的地基：\textbf{极限}。
